\documentclass[11pt]{article}
\usepackage[T1]{fontenc}
\usepackage{lmodern}
\usepackage[margin=1in]{geometry}
\usepackage{microtype}
\usepackage{amsmath,amssymb,amsthm,mathtools}
\usepackage[colorlinks=true,linkcolor=blue,citecolor=blue,urlcolor=blue]{hyperref}
\hypersetup{pdftitle={The Matrix Multiplication Bound ω ≤ 9/4 over Arbitrary Fields},pdfauthor={}}
\newtheorem{theorem}{Theorem}[section]
\newtheorem{proposition}[theorem]{Proposition}
\newtheorem{lemma}[theorem]{Lemma}
\newtheorem{corollary}[theorem]{Corollary}
\theoremstyle{remark}
\newtheorem{remark}[theorem]{Remark}
\DeclareMathOperator{\rank}{R}
\newcommand{\C}{\mathbb C}
\newcommand{\Q}{\mathbb Q}
\newcommand{\N}{\mathbb N}
\newcommand{\R}{\mathbb R}
\newcommand{\eps}{\varepsilon}
\newcommand{\degen}{\mathrel{\rightsquigarrow}}
\numberwithin{equation}{section}
\title{The Matrix Multiplication Bound\\$\omega\leq9/4$ over Arbitrary Fields}
\author{}
\date{}
\begin{document}
\maketitle

\begin{abstract}
OpenAI proved that the matrix multiplication exponent over the complex
numbers is at most $9/4$. We extend that argument to every field. The
extension replaces a Fourier period that can vanish in positive
characteristic, removes an integer-node choice from the Lean interpolation
proof, and descends exact tensor decompositions from an algebraic closure
with a fixed overhead. We describe these changes and their formalization,
taking the underlying $9/4$ construction from OpenAI's proof.
\end{abstract}

\section{Statement and relation to the original result}

For a field $F$, write $\omega(F)$ for the infimum of the real numbers
$\tau$ such that, for every $\eps>0$, square matrices of size $n$ over $F$
can be multiplied by division-free arithmetic circuits of cost
$O_{F,\eps}(n^{\tau+\eps})$. Scalar additions, subtractions, and
multiplications each cost one; inputs and constant loads are free.
Circuits and their constants may depend on the size, and the constant
implicit in the bound may depend on $F$ and $\eps$.

\begin{theorem}\label{thm:main}
For every field $F$ and every $\eps>0$, there is a constant $C>0$ such
that, for every integer $n\geq1$, there is a division-free arithmetic
circuit over $F$ computing $n\times n$ matrix multiplication with at
most $Cn^{9/4+\eps}$ operations. In particular,
\[
\omega(F)\leq\frac94.
\]
The bound is obtained from exact identities of tensor coefficients.
\end{theorem}

The original bound over $\C$, its auxiliary-separation construction, and
its spectral and polynomial-growth argument are due to OpenAI
\cite{OAI}. This note concerns the additional field scope. We use the
manuscript and Lean sources at the pinned revision in \cite{OAI,OAILean};
the formal extension is archived in \cite{Formalization}.

The characteristic-zero consequence is already implicit in
\cite[Remark~5.2]{OAI}. For a fixed finite rank decomposition, its
coefficient equations have rational coefficients. A complex solution
therefore gives a solution over $\overline{\Q}$, and the finitely many
coefficients lie in a number field. Coordinate simulation over $\Q$
has constant overhead once the decomposition is fixed. The resulting
rational coefficient identities extend to every characteristic-zero
field. The extension degree may depend on $\eps$, which is permitted
in the definition of the exponent.

This reasoning does not by itself cover every positive characteristic:
reducing a rational scheme modulo a prime can fail at its denominators,
and the exceptional primes can depend on the scheme chosen for $\eps$.
Instead, for each field $F$ we work directly over $K=\overline F$.
Section~\ref{sec:closed} explains the modifications over $K$;
Section~\ref{sec:descent} descends the result to $F$.

We distinguish the circuit exponent from the exact coefficient-rank
exponent. Let
\[
T_n=\sum_{i,j,k=1}^{n}x_{ij}y_{jk}z_{ki},\qquad
\nu(F)=\inf_{n\geq2}\log_n\rank_F(T_n),
\]
where $\rank_F$ is tensor rank over $F$. The proof gives
$\nu(F)\leq9/4$ and then $\omega(F)\leq\nu(F)$. For unrestricted
arithmetic circuits over finite fields, functional correctness alone
need not imply a formal polynomial identity; the tensor argument proves
exact coefficient identities.

\section{The construction in positive characteristic}\label{sec:closed}

Fix an algebraically closed field $K$. Such a field is infinite,
regardless of its characteristic. We retain the tensor restrictions,
degenerations, and real-valued tensor characters of \cite{OAI}.

\subsection{A nonvanishing Fourier period}

The finite separation construction in \cite[Proposition~3.1]{OAI}
uses a Fourier period $5M$ for $M\geq1$ blocks. Normalization divides
by this period. In characteristic $p$ this is invalid when $p\mid5M$;
in that case there is also no primitive $(5M)$-th root of unity.
Replace the period by
\begin{equation}\label{eq:period}
L=L_K(M)=
\begin{cases}
5M,&(5M)1_K\ne0,\\
5M+1,&(5M)1_K=0.
\end{cases}
\end{equation}
Then $5M\leq L\leq6M$ and $L1_K\ne0$, since consecutive integers
cannot both have zero image in a field. The polynomial $X^L-1$ splits
with distinct roots in $K$. Its roots form a finite subgroup of
$K^\times$, hence a cyclic group, so $K$ contains a primitive
$L$-th root $\zeta$.

\begin{lemma}\label{lem:fourier}
For $e\in\mathbb Z$,
\[
\frac1L\sum_{r=0}^{L-1}\zeta^{re}
=\begin{cases}1,&L\mid e,\\0,&L\nmid e.\end{cases}
\]
If $|e|<L$, the value is $1$ precisely when $e=0$.
\end{lemma}
\begin{proof}
If $L\mid e$, all summands equal $1$. Otherwise $\zeta^e\ne1$,
and the geometric-sum identity gives zero. Both divisions are valid:
$L1_K$ and $\zeta^e-1$ are nonzero in their respective cases.
The last statement is an assertion about divisibility of integers.
\end{proof}

In the substitutions of \cite[Proposition~3.1]{OAI}, the Fourier
phase is $e=u-v+2(g-h)$, with all four labels in $\{1,\ldots,M\}$.
It satisfies $|e|\leq3(M-1)<L$. Thus Lemma~\ref{lem:fourier}
retains exactly the same terms as the original filter. The unchanged
weight calculation assigns total weight $(g-h)^2$ to a surviving
term; its zero-weight part forces $g=h$ and $u=v$. These weights are
integer exponents of a formal parameter, so they are not reduced
modulo the characteristic. In particular the argument also works
in characteristic two.

The resulting separation uses $L$ copies of the source tensor instead
of $5M$ copies. This changes a bounded finite overhead, not the
entropy inequality. Indeed, in the type-counting step of
\cite[Corollary~3.2]{OAI}, let $M_N$ be the number of words of a
fixed rational type $q$, where $N$ runs through multiples of a common
denominator of $q$, and put
$C_N=\prod_i\lambda(A_i)^{Nq_i}$. Here the $A_i$ are the sectors
of $A$, $\lambda$ is a tensor character, and
$\lambda(B_X(m))=m^{p_X}$ for the oriented dot-product tensor.
Separation and degeneration monotonicity give
\[
6M_N\lambda(A)^N\ \geq\ L_K(M_N)\lambda(A)^N
\ \geq\ M_N^{1+p_X}C_N.
\]
After canceling $M_N$ and taking $N$-th roots, the extra factor is
$6^{1/N}\to1$. The same entropy limit and continuity argument
therefore give the original inequality. All quantities in this
limit are real numbers, not elements of $K$.

\subsection{Interpolation over the algebraic closure}

The manuscript's interpolation argument already allows arbitrary
distinct nonzero nodes \cite[Lemma~2.3]{OAI}. The corresponding
original Lean implementation chooses the images of
$1,2,\ldots,D+1$ in $\C$. That implementation choice must be
replaced in positive characteristic: integer images can coincide
or vanish. An infinite field supplies any finite number of distinct
nonzero nodes, with invertible pairwise differences.

For clarity, the part of interpolation used here can be expressed
without characteristic assumptions as follows. Suppose a polynomial
tensor family $G(t)$ has degree at most $D$, has $G(0)=B$, and every
$G(t)$ with $t\ne0$ is a restriction of $A$. For each $m\geq1$,
the family $G(t)^{\otimes m}$ has degree at most $Dm$ and constant
coefficient $B^{\otimes m}$. Interpolating at $Dm+1$ distinct nonzero
nodes expresses that coefficient as a linear combination of the
evaluations. Hence
\[
\lambda(B)^m\leq(Dm+1)\lambda(A)^m.
\]
Taking $m$-th roots proves $\lambda(B)\leq\lambda(A)$.
If a degeneration initially has a nonzero leading degree, divide
out that power of the parameter before applying this argument;
the nonzero-node condition permits this normalization.
The same choice of nodes gives the usual evaluation and interpolation
bound for polynomial multiplication over $K$.

\subsection{The retained 9/4 argument}

With these two choices made, the remaining proof of \cite{OAI}
applies over $K$. The determinant filtration uses polynomial-domain
identities and bases with unit pivots; it does not require division
by the characteristic or a semisimple splitting. The sector
construction uses integer indices and weights. The real growth
inequalities in \cite[Sections~4--5]{OAI} have no dependence on the
characteristic of the coefficient field. Finally, the separation and
fixed-point arguments in \cite[Appendix~A]{OAI} act on real-valued
functions on the tensor semiring. They do not impose a topology on
$K$. A natural number in that semiring counts direct summands and
is not its image in $K$.

Consequently the detecting-character argument, with the same
polynomial-growth estimate, gives
\begin{equation}\label{eq:closed}
\nu(K)\leq\frac94\qquad(K\text{ algebraically closed}).
\end{equation}
This step uses the original proof rather than a new derivation of
the numerical bound. In the formalization, its endpoint is
\nolinkurl{AuxiliarySeparation.exactRankExponent_le_nine_quarters}.

\section{Descent with a fixed coefficient algebra}\label{sec:descent}

Field-extension invariance is classical; see Sch\"onhage's rank-descent
argument and characteristic-invariance theorem
\cite[Theorems~2.7--2.8]{Schonhage}. The descent principle below
isolates the order of quantifiers used in our formalization. It uses
finite-dimensional linear algebra and applies without a separability
assumption.

\begin{lemma}\label{lem:projection}
Let $S$ be a nonzero commutative $F$-algebra of finite dimension $d$.
If a finite coefficient tensor $T$ over $F$ has an $r$-term rank-one
decomposition after scalar extension to $S$, then
\[
\rank_F(T)\leq d^2r.
\]
\end{lemma}
\begin{proof}
Choose an $F$-basis $e_1,\ldots,e_d$ of $S$ and an $F$-linear map
$\pi:S\to F$ with $\pi(1)=1$. Such a map exists because $1\ne0$;
extend $1$ to a basis and take its coordinate functional. Thus
$\pi(c1_S)=c$ for $c\in F$.

Write the decomposition as
$T(x,y,z)1_S=\sum_{\ell=1}^r a_\ell(x)b_\ell(y)c_\ell(z)$.
Expand the first two factors in the chosen basis:
$a_\ell(x)=\sum_i a_{\ell i}(x)e_i$ and
$b_\ell(y)=\sum_j b_{\ell j}(y)e_j$, with coefficients in $F$.
Applying $\pi$ gives the coefficient identity
\[
T(x,y,z)=\sum_{\ell=1}^r\sum_{i,j=1}^d
a_{\ell i}(x)b_{\ell j}(y)\,
\pi\bigl(e_i e_j c_\ell(z)\bigr).
\]
Each summand separates the three coordinates and has coefficients
in $F$. There are at most $rd^2$ summands. No multiplicativity of
$\pi$ or field trace is used.
\end{proof}

\begin{proposition}\label{prop:powers}
Let $E/F$ be an algebraic field extension, and let $T$ be a finite
coefficient tensor over $F$. If $\rank_E(T_E)\leq r$, then there
exists an integer $d\geq1$ such that, for every $m\geq1$,
\begin{equation}\label{eq:descent}
\rank_F(T^{\otimes m})\leq d^2r^m.
\end{equation}
\end{proposition}
\begin{proof}
Fix one $r$-term decomposition over $E$. Its finitely many coefficients
generate an $F$-subalgebra $S\subseteq E$ of finite dimension $d$,
because each generator is algebraic over $F$. The chosen decomposition
already exists over $S$, and taking its $m$-th tensor power uses only
products of coefficients in the same algebra $S$. It therefore has
at most $r^m$ terms over $S$. Apply Lemma~\ref{lem:projection} to
this powered decomposition.
\end{proof}

The order is $\exists d\,\forall m$: the coefficient algebra is
fixed before the tensor power varies. Descending one block first
and then powering its descended decomposition would instead give
the weaker bound $(d^2r)^m$. Equation~\eqref{eq:descent} avoids
charging the extension cost at each recursive level. The proof
works for inseparable extensions and for imperfect base fields.

\begin{corollary}\label{cor:exponent}
For every algebraic field extension $E/F$,
$\nu(F)\leq\nu(E)$. In particular $\nu(F)\leq\nu(\overline F)$.
\end{corollary}
\begin{proof}
Fix $u\geq2$ and a decomposition of $T_u$ over $E$ of length
$r=\rank_E(T_u)$. Since $T_u^{\otimes m}$ is $T_{u^m}$ after
reindexing, Proposition~\ref{prop:powers} yields, for every $m\geq1$,
\[
\nu(F)\leq\log_{u^m}\rank_F(T_{u^m})
\leq\log_u r+\frac{2\log d}{m\log u}.
\]
Let $m\to\infty$, then take the infimum over $u\geq2$.
This does not require the infimum defining $\nu(E)$ to be attained.
\end{proof}

\section{Arithmetic circuits and the formal statement}

\begin{proof}[Proof of Theorem~\ref{thm:main}]
Apply \eqref{eq:closed} to $K=\overline F$ and then
Corollary~\ref{cor:exponent} to obtain $\nu(F)\leq9/4$.
Fix $0<\delta<\eps$. The definition of the infimum supplies a
block size $u\geq2$ and an exact rank decomposition over $F$ of
length $r<u^{9/4+\delta}$.

The generic block-program construction from the original Lean
development \cite{OAILean} implements this bilinear formula
recursively, including the scalar operations in its linear
combinations. On powers of $u$, its cost satisfies
\[
S(uv)\leq rS(v)+6ru^2v^2,\qquad S(1)=1.
\]
The flattening bound gives $r\geq u^2$. Thus the cost is
$O(N^{\log_u r})$ when $r>u^2$, and $O(N^2\log N)$ when $r=u^2$.
Padding to a power of the fixed $u$ and using the positive exponent
slack gives $O_{F,\eps}(n^{9/4+\eps})$ for every $n\geq1$.
\end{proof}

The archived formalization \cite{Formalization}, written in Lean~4
\cite{Lean} using Mathlib \cite{Mathlib}, exports the theorem
in namespace \texttt{OAI.MatrixMultiplication}. With the real type
written as \texttt{Real}, its statement is
\begin{verbatim}
theorem omega_le_nine_quarters (F : Type*) [Field F] :
    Arithmetic.omega F <= (9 : Real) / 4
\end{verbatim}
The original statement substitutes $\C$ for $F$. The definitions
of the arithmetic model are preserved; the current source includes
only a mechanical module-system port of the protected upstream
specification. The original complex conclusion remains a specialization.

Two supporting results make the scope explicit. The exact-rank theorem
proves $\nu(F)\leq9/4$ using equality at every tensor coordinate.
The arithmetic theorem has quantifiers
\[
\forall\eps>0\ \exists C>0\ \forall n\geq1\ \exists P:
\quad P\text{ correctly multiplies over }F,\qquad
\operatorname{cost}(P)\leq Cn^{9/4+\eps}.
\]
Here $P$ must work for every input pair of its size. Over finite
fields, its correctness predicate is functional; the separate exact
coefficient theorem supplies the stronger algebraic guarantee.
The formal result is a family of arithmetic circuits. It does not
assert an $O(n^{9/4})$ endpoint bound, effective uniform generation,
or a bit-complexity bound.

\begin{flushleft}
\small
The source correspondence is short. Within the matrix-multiplication
subtree of \cite{Formalization}, the Fourier changes are in
\nolinkurl{AuxiliarySeparation/Separation}; interpolation is in
\nolinkurl{AuxiliarySeparation/Polynomial/Interpolation.lean}.
The two descent results are implemented in
\nolinkurl{Arithmetic/FieldDescent.lean}. The exponent comparison is in
\nolinkurl{AuxiliarySeparation/Arithmetic/FieldExtension.lean}.
The public entry point is \nolinkurl{AllFields.lean}.
\end{flushleft}

The versioned archive records the proof snapshot and its mechanical
verification. Its six Comparator claims include the arithmetic bound,
infimum side conditions, and explicit exact coefficient identities.
The recorded Palomar run checks those claims using Lean and the
independent kernels listed in its report. The allowed axioms are
\texttt{propext}, \texttt{Classical.choice}, and \texttt{Quot.sound}.
These are verification claims about the recorded source snapshot,
not a claim of human peer review or a check of this prose document.

\paragraph{Attribution and development.}
OpenAI supplied the original $9/4$ construction and complex-field proof
\cite{OAI,OAILean}. The contribution described here is its adaptation
to arbitrary positive characteristic, the explicit descent argument,
and their formal verification. No historical-priority claim is made
for the scalar-descent principle. The extension and Lean formalization
were developed with consumer-accessible GPT-6 Astra and GPT-6.1 Sol
through Codex, under human direction. The repository records that
provenance. This note was prepared with AI assistance.

\begin{thebibliography}{9}
\bibitem{OAI}
OpenAI, \emph{An Upper Bound of $9/4$ for the Matrix Multiplication
Exponent}, October 2, 2026.
\href{https://github.com/openai/math/blob/adc7f1241b42e322a6451854ab7e4b4c146bf78a/preprints/Matrix-Multiplication-Nine-Fourths-October-2-2026/paper.pdf}{Preprint},
\href{https://github.com/openai/math/blob/adc7f1241b42e322a6451854ab7e4b4c146bf78a/preprints/Matrix-Multiplication-Nine-Fourths-October-2-2026/build/paper.tex}{manuscript source}.
Revision \texttt{adc7f1241b42e322a6451854ab7e4b4c146bf78a}.

\bibitem{OAILean}
OpenAI, \emph{Lean sources for matrix multiplication}, 2026.
\href{https://github.com/openai/math/tree/adc7f1241b42e322a6451854ab7e4b4c146bf78a/lean/OAI/LinearAlgebra/MatrixMultiplication}{MatrixMultiplication subtree},
revision \texttt{adc7f1241b42e322a6451854ab7e4b4c146bf78a}.

\bibitem{Formalization}
S.~Navot, \emph{Matrix multiplication exponent at most $9/4$ over
arbitrary fields: a Lean formalization}, version 1.0.0, October 7, 2026.
\href{https://doi.org/10.5281/zenodo.23219128}{doi:10.5281/zenodo.23219128}.
\href{https://github.com/selanavot/matrix-multiplication-all-fields/tree/c88bb830aeb6ebd83f1e48a7bdae566ff5e0924d}{Archived source revision}
\texttt{c88bb830aeb6ebd83f1e48a7bdae566ff5e0924d};
verification provenance is in \texttt{docs/releases/v1.0.0.md}.

\bibitem{Schonhage}
A.~Sch\"onhage, \emph{Partial and total matrix multiplication},
SIAM Journal on Computing \textbf{10}(3) (1981), 434--455.
\href{https://doi.org/10.1137/0210032}{doi:10.1137/0210032}.

\bibitem{Lean}
L.~de Moura and S.~Ullrich, \emph{The Lean 4 theorem prover and
programming language}, in Automated Deduction---CADE~28,
Lecture Notes in Computer Science \textbf{12699} (2021), 625--635.
\href{https://doi.org/10.1007/978-3-030-79876-5_37}{doi:10.1007/978-3-030-79876-5\_37}.

\bibitem{Mathlib}
The mathlib Community, \emph{The Lean mathematical library},
Proceedings of the 9th ACM SIGPLAN International Conference on
Certified Programs and Proofs (CPP~2020), 367--381.
\href{https://doi.org/10.1145/3372885.3373824}{doi:10.1145/3372885.3373824}.
\end{thebibliography}
\end{document}
